\documentclass[11pt,a4paper]{article}
\usepackage[utf8]{inputenc}
\usepackage[T1]{fontenc}
\usepackage{amsmath,amssymb,mathtools}
\usepackage{geometry}
\geometry{margin=2.6cm}
\usepackage{xcolor,mdframed,enumitem}
\usepackage[expansion=false]{microtype}
\usepackage[bookmarks=false]{hyperref}
\hypersetup{colorlinks=true,linkcolor=blue!55!black,hypertexnames=false}
\definecolor{blueT}{RGB}{32,92,155}
\definecolor{orangeL}{RGB}{205,105,20}
\definecolor{greenP}{RGB}{34,125,72}
\definecolor{violetR}{RGB}{110,62,145}
\definecolor{redK}{RGB}{175,40,40}
\newmdenv[backgroundcolor=blue!3,linecolor=blueT,linewidth=1.2pt,roundcorner=4pt]{theorembox}
\newmdenv[backgroundcolor=orange!4,linecolor=orangeL,linewidth=1.2pt,roundcorner=4pt]{lemmabox}
\newmdenv[backgroundcolor=green!3,linecolor=greenP,linewidth=1.2pt,roundcorner=4pt]{proofbox}
\newmdenv[backgroundcolor=violet!3,linecolor=violetR,linewidth=1.2pt,roundcorner=4pt]{remarkbox}
\newmdenv[backgroundcolor=red!3,linecolor=redK,linewidth=1.4pt,roundcorner=4pt]{resultbox}
\newcommand{\N}{\mathbb N}
\newcommand{\R}{\mathbb R}
\newcommand{\eps}{\varepsilon}
\newcommand{\e}{\mathrm e}
\title{Jean Dieudonn\'e -- Calcul infinit\'esimal\\
\large R\'esolution rigoureuse des exercices 25 et 26, page 119}
\author{Mode chercheur senior}
\date{11 juillet 2026}
\begin{document}
\maketitle
\tableofcontents

\begin{remarkbox}
\textbf{Statut \'epist\'emique.} Les deux r\'esultats sont inconditionnels.
Les preuves suivent les d\'ecoupages indiqu\'es par Dieudonn\'e et n'utilisent
pas de th\'eor\`eme de convergence domin\'ee pour les s\'eries discr\`etes.
Chaque estimation de queue est quantitative.
\end{remarkbox}

\section{Exercice 25 : somme de grandes puissances}

\subsection{\'Enonc\'e reformul\'e}
\begin{theorembox}
Soit $\alpha>0$. D\'emontrer que
\[
 1^{\alpha n}+2^{\alpha n}+\cdots+n^{\alpha n}
 \sim \frac{n^{\alpha n}}{1-\e^{-\alpha}}
 \qquad(n\to+\infty).
\]
L'indication propose de majorer
$\sum_{k=1}^{n-1}|(1-k/n)^{\alpha n}-\e^{-\alpha k}|$ en la coupant \`a
$k=\lfloor n^\lambda\rfloor$, avec $\lambda$ convenable.
\end{theorembox}

\subsection{Conventions et pr\'erequis}
En posant $k=n-m$, on obtient l'identit\'e exacte
\begin{equation}
 \frac1{n^{\alpha n}}\sum_{m=1}^n m^{\alpha n}
 =\sum_{k=0}^{n-1}\left(1-\frac{k}{n}\right)^{\alpha n}. \tag{1}
\end{equation}
Pour tout $x\in[0,1[$,
\begin{equation}
 \log(1-x)\leq-x. \tag{2}
\end{equation}
Pour $x\in[0,1/2]$, la s\'erie de $-\log(1-x)$ donne
\begin{equation}
 0\leq-\log(1-x)-x=\sum_{r\geq2}\frac{x^r}{r}
 \leq\frac{x^2}{2(1-x)}\leq x^2. \tag{3}
\end{equation}

\subsection{Lemmes interm\'ediaires}
\begin{lemmabox}
\textbf{Lemme 1.} Pour $0\leq k\leq n/2$,
\[
 0\leq \e^{-\alpha k}-\left(1-\frac{k}{n}\right)^{\alpha n}
 \leq \frac{\alpha k^2}{n}\e^{-\alpha k}.
\]
\end{lemmabox}
\begin{proofbox}
Posons
\[
 u_{n,k}:=\alpha n\left(-\log\left(1-\frac{k}{n}\right)-\frac{k}{n}\right).
\]
Les in\'egalit\'es (2)--(3) donnent $0\leq u_{n,k}\leq\alpha k^2/n$ et
\[
 \left(1-\frac{k}{n}\right)^{\alpha n}=\e^{-\alpha k}\e^{-u_{n,k}}.
\]
Comme $0\leq1-\e^{-u}\leq u$ pour $u\geq0$, la conclusion suit.
\end{proofbox}

\begin{lemmabox}
\textbf{Lemme 2.} Pour toute suite d'entiers $K_n\to\infty$ telle que
$K_n\leq n/2$ \`a partir d'un certain rang,
\[
 \sum_{k=0}^{K_n}\left|\left(1-\frac{k}{n}\right)^{\alpha n}
 -\e^{-\alpha k}\right|=O_\alpha(n^{-1}).
\]
\end{lemmabox}
\begin{proofbox}
Par le Lemme 1,
\[
 \sum_{k=0}^{K_n}\left|\left(1-\frac{k}{n}\right)^{\alpha n}
 -\e^{-\alpha k}\right|
 \leq\frac{\alpha}{n}\sum_{k=0}^{\infty}k^2\e^{-\alpha k}.
\]
La s\'erie du membre droit converge, puisque le quotient de deux termes
cons\'ecutifs tend vers $\e^{-\alpha}<1$.
\end{proofbox}

\subsection{Preuve principale par le d\'ecoupage $n^\lambda$}
Fixons une fois pour toutes $\lambda\in(0,1)$, par exemple $\lambda=1/3$, et
posons $K_n=\lfloor n^\lambda\rfloor$. Pour $n$ assez grand, $K_n\leq n/2$.
Notons
\[
 D_n:=\sum_{k=0}^{n-1}\left(1-\frac{k}{n}\right)^{\alpha n}
       -\sum_{k=0}^{\infty}\e^{-\alpha k}.
\]
Nous d\'ecomposons $D_n$ exactement en trois morceaux :
\begin{align*}
D_n={}&\sum_{k=0}^{K_n}
 \left\{\left(1-\frac{k}{n}\right)^{\alpha n}-\e^{-\alpha k}\right\}\\
&+\sum_{k=K_n+1}^{n-1}\left(1-\frac{k}{n}\right)^{\alpha n}
-\sum_{k=K_n+1}^{\infty}\e^{-\alpha k}.
\end{align*}
Le premier morceau est $O_\alpha(n^{-1})$ par le Lemme 2. D'apr\`es (2),
\[
 0\leq\left(1-\frac{k}{n}\right)^{\alpha n}\leq\e^{-\alpha k},
 \qquad 0\leq k<n.
\]
Par cons\'equent, chacun des deux morceaux de queue est major\'e en valeur
absolue par
\[
 \sum_{k=K_n+1}^{\infty}\e^{-\alpha k}
 =\frac{\e^{-\alpha(K_n+1)}}{1-\e^{-\alpha}}\longrightarrow0.
\]
Ainsi $D_n\to0$, et la somme g\'eom\'etrique fournit
\[
 \sum_{k=0}^{n-1}\left(1-\frac{k}{n}\right)^{\alpha n}
 \longrightarrow\sum_{k=0}^{\infty}\e^{-\alpha k}
 =\frac1{1-\e^{-\alpha}}.
\]
En revenant \`a (1), nous obtenons
\begin{resultbox}
\[
 \boxed{\displaystyle
 1^{\alpha n}+2^{\alpha n}+\cdots+n^{\alpha n}
 \sim\frac{n^{\alpha n}}{1-\e^{-\alpha}}.}
\]
\end{resultbox}

\subsection{Remarque sur le choix de $\lambda$}
\begin{remarkbox}
Le d\'ecoupage isole la couche limite $n-m=O(n^\lambda)$, seule zone o\`u les
termes restent visibles apr\`es division par $n^{\alpha n}$. Avec l'estimation
quantitative du Lemme 1, tout $\lambda\in(0,1)$ convient. Si l'on se contente
de l'approximation uniforme plus grossi\`ere
$\alpha n\log(1-k/n)=-\alpha k+O(n^{2\lambda-1})$ puis que l'on somme sans
tenir compte du facteur $\e^{-\alpha k}$, il faut imposer une condition plus
restrictive. Le Lemme 1 \'evite cette perte.
\end{remarkbox}

\section{Exercice 26 : somme de puissances de factorielles}

\subsection{\'Enonc\'e reformul\'e}
\begin{theorembox}
Soit $\alpha>0$. D\'emontrer que
\[
 \sum_{p=1}^{n}(p!)^{-\alpha/n}
 \sim\frac1\alpha\frac{n}{\log n}
 \qquad(n\to+\infty).
\]
L'indication demande de comparer $(p!)^{-\alpha/n}$ \`a $p^{-\alpha p/n}$ et
de couper la somme \`a $p=\lfloor\eps n\rfloor$, o\`u $\eps>0$ est arbitraire.
\end{theorembox}

\subsection{Conventions et estimations pr\'eliminaires}
Posons
\[
 N_n:=\frac{n}{\log n},\qquad
 S_n:=\sum_{p=1}^{n}(p!)^{-\alpha/n},\qquad
 T_n:=\sum_{p=1}^{n}p^{-\alpha p/n}.
\]
Les in\'egalit\'es int\'egrales appliqu\'ees \`a $\log$ donnent, pour $p\geq1$,
\begin{equation}
 p\log p-p+1\leq\log(p!)\leq p\log p. \tag{4}
\end{equation}
Il en r\'esulte
\begin{equation}
 p^{-\alpha p/n}leq(p!)^{-\alpha/n}
 \leq \e^{\alpha(p-1)/n}p^{-\alpha p/n}. \tag{5}
\end{equation}

\subsection{Lemme central : somme mod\`ele}
\begin{lemmabox}
\textbf{Lemme 3.} On a
\[
 T_n\sim\frac1\alpha N_n.
\]
\end{lemmabox}
\begin{proofbox}
Fixons $0<\delta<M$. Pour $p$ tel que $\delta N_n\leq p\leq MN_n$, posons
$u_p=p/N_n$. Alors, uniform\'ement pour $u_p\in[\delta,M]$,
\begin{align*}
\frac{\alpha p\log p}{n}
&=\alpha u_p\frac{\log(u_p n/\log n)}{\log n}\\
&=\alpha u_p\left(1-\frac{\log\log n}{\log n}
  +\frac{\log u_p}{\log n}\right)
=\alpha u_p+o(1).
\end{align*}
Ainsi
\[
 \frac1{N_n}\sum_{\delta N_n\leq p\leq MN_n}p^{-\alpha p/n}
 \longrightarrow\int_\delta^M\e^{-\alpha u}\,du, \tag{6}
\]
par la d\'efinition des sommes de Riemann.

Il reste \`a contr\^oler les deux extr\'emit\'es. Pour $p<\delta N_n$, chaque
terme est au plus $1$, donc la contribution divis\'ee par $N_n$ est au plus
$\delta+o(1)$. Pour la queue, s\'eparons d'abord $p<\sqrt n$, qui ne peut
intervenir que si $M N_n<\sqrt n$; pour $n$ assez grand et $M$ fix\'e, on a au
contraire $MN_n>\sqrt n$. Si $p\geq\sqrt n$, alors
$\log p\geq\frac12\log n$, donc
\[
 p^{-\alpha p/n}\leq
 \exp\left(-\frac{\alpha p\log n}{2n}\right)
 =\exp\left(-\frac{\alpha}{2}\frac{p}{N_n}\right).
\]
Par comparaison avec une s\'erie g\'eom\'etrique,
\begin{align*}
\frac1{N_n}\sum_{p>MN_n}p^{-\alpha p/n}
&\leq\frac1{N_n}\sum_{p>MN_n}\e^{-\alpha p/(2N_n)}\\
&\leq C_\alpha\e^{-\alpha M/2}
\end{align*}
pour $n$ assez grand, avec une constante ind\'ependante de $M$ et $n$.
Faisons successivement $n\to\infty$, $\delta\downarrow0$, puis $M\to\infty$
dans (6). Nous obtenons
\[
 \frac{T_n}{N_n}\longrightarrow\int_0^\infty\e^{-\alpha u}\,du=\frac1\alpha.
\]
\end{proofbox}

\subsection{Preuve principale avec la coupure $\lfloor\eps n\rfloor$}
Fixons $\eps\in(0,1)$. D'apr\`es (5), pour $p\leq\eps n$,
\[
 p^{-\alpha p/n}\leq(p!)^{-\alpha/n}
 \leq\e^{\alpha\eps}p^{-\alpha p/n}. \tag{7}
\]
Nous montrons que la partie $p>\eps n$ est n\'egligeable \`a l'\'echelle $N_n$.
Pour $n$ assez grand et $p>\eps n$, l'in\'egalit\'e (4) entra\^ine
\[
 (p!)^{-\alpha/n}
 \leq\exp\left[-\frac{\alpha p}{n}(\log p-1)\right]
 \leq\exp\left[-\frac{\alpha p\log n}{2n}\right]
 =\e^{-\alpha p/(2N_n)}.
\]
Donc
\begin{equation}
 0\leq\frac1{N_n}\sum_{p>\eps n}^{n}(p!)^{-\alpha/n}
 \leq C_\alpha\exp\left(-\frac{\alpha\eps}{2}\log n\right)
 \longrightarrow0. \tag{8}
\end{equation}
La m\^eme estimation vaut pour la queue de $T_n$. En sommant (7), puis en
utilisant (8) et le Lemme 3, nous obtenons
\[
 \frac1\alpha
 \leq\liminf_{n\to\infty}\frac{S_n}{N_n}
 \leq\limsup_{n\to\infty}\frac{S_n}{N_n}
 \leq\frac{\e^{\alpha\eps}}\alpha.
\]
Cette in\'egalit\'e est valable pour tout $\eps>0$. En faisant tendre
$\eps\downarrow0$, les membres extr\^emes ont la m\^eme limite $1/\alpha$.
Ainsi
\begin{resultbox}
\[
 \boxed{\displaystyle
 (1!)^{-\alpha/n}+(2!)^{-\alpha/n}+\cdots+(n!)^{-\alpha/n}
 \sim\frac1\alpha\frac{n}{\log n}.}
\]
\end{resultbox}

\subsection{Interpr\'etation asymptotique}
\begin{remarkbox}
L'\'echelle naturelle est $p=u,n/\log n$. En effet,
\[
 \frac{\alpha}{n}\log(p!)\sim\frac{\alpha p\log p}{n}\sim\alpha u,
\]
de sorte que le terme g\'en\'eral ressemble \`a $\e^{-\alpha u}$ et que le pas
en $u$ vaut $1/N_n$. Le coefficient $1/\alpha$ est donc l'int\'egrale
$\int_0^\infty\e^{-\alpha u}\,du$.
\end{remarkbox}

\section{Auto-v\'erification et r\'ef\'erence}
\begin{itemize}[leftmargin=*]
\item Les hypoth\`eses $\alpha>0$ sont utilis\'ees pour la convergence des
queues g\'eom\'etriques.
\item Aucun passage somme--limite infini n'est laiss\'e implicite : les zones
centrales sont trait\'ees par sommes de Riemann et les queues sont major\'ees.
\item Les coupures $n^\lambda$ et $\lfloor\eps n\rfloor$ sont celles sugg\'er\'ees
dans l'ouvrage.
\end{itemize}
\medskip
J. Dieudonn\'e, \emph{Calcul infinit\'esimal}, chapitre III,
``D\'eveloppements asymptotiques'', exercices 25 et 26, p.~119.

\end{document}
